\documentclass[10pt,a4paper]{amsart}
\usepackage[margin=19mm]{geometry}
\usepackage{amsmath,amssymb,mathtools}
\usepackage[scr=rsfs,cal=euler]{mathalfa}
\usepackage{xfrac}
\usepackage[unicode,hidelinks,bookmarks=false]{hyperref}
\newcommand{\ie}{i.e.\ }
\newcommand{\cf}{cf.\ }
\newcommand{\resp}{respectively\ }
\newcommand{\Subsection}{\subsection}
\providecommand{\nobreakdash}{\nobreak\hbox{-}}
\setlength{\emergencystretch}{2em}
\multlinegap0pt
\usepackage{bm}
\usepackage{bbm}
\usepackage{dsfont}
\usepackage{xspace}
\usepackage{cancel}
\usepackage{mathtools}
\usepackage{amsthm}
\makeatletter
\@ifundefined{thm}{\newtheorem{thm}{Thm}}{}
\@ifundefined{claim}{\newtheorem{claim}[thm]{\bf Claim}}{}
\@ifundefined{cor}{\newtheorem{cor}{\bf Corollary}}{}
\@ifundefined{lem}{\newtheorem{lem}[thm]{Lemma}}{}
\makeatother
\newcommand{\Z}{\mathbb{Z}}
\newcommand{\barB}{\overline{B}}
\newcommand{\bfa}{\mathbf{a}}
\newcommand{\diffSj}{\Delta {S}^{(j)}}
\newcommand{\dist}{{\rm dist}}
\newcommand{\length}{{\rm length}}
\newcommand{\mcD}{\mathcal{D}}
\newcommand{\mcE}{\mathcal{E}}
\newcommand{\mcG}{\mathcal{G}}
\newcommand{\mcI}{\mathcal{I}}
\newcommand{\mcR}{\mathcal{R}}
\newcommand{\partialSjminusone}{\partial {S}^{(j-1)}}
\newcommand{\partialSj}{\partial {S}^{(j)}}
\newcommand{\smallSjminusone}{s^{(j-1)}}
\newcommand{\smallSj}{s^{(j)}}
\newcommand{\spec}{{\rm Spec}}
\newcommand{\tm}{\widetilde{m}}
\newcommand{\tq}{\tilde{q}}
\newcommand{\tr}{\tilde{r}}
\newcommand{\ts}{\tilde{s}}
\newcommand{\whn}{\widehat{n}}
\newcommand{\wtGjminusone}{\widetilde{G}^{(j-1)}}
\newcommand{\wtSjminusone}{\widetilde{S}^{(j-1)}}
\newcommand{\wtSj}{\widetilde{S}^{(j)}}
\renewcommand{\tt}{\tilde{t}}
\title[Anderson Localization for the hierarchical Anderson-Bernoull]{Anderson Localization for the hierarchical Anderson-Bernoulli model on $\mathbb{Z}^d$}
\author{Shihe Liu, Yunfeng Shi, Zhifei Zhang}
\date{Source: arXiv:2604.18989v1 (CC BY 4.0)}
\begin{document}
\maketitle

\noindent\emph{Source and licence.} Anderson Localization for the hierarchical Anderson-Bernoulli model on $\mathbb{Z}^d$. Version arXiv:2604.18989v1 (CC BY 4.0), distributed under the Creative Commons Attribution 4.0 International licence, as recorded in the arXiv licence metadata for this exact version and shown on the abstract page https://arxiv.org/abs/2604.18989v1. Full rights notice: licenses/arxiv-2604.18989-CC-BY-4.0.txt.\par\smallskip

\noindent\textit{Editorial scope.} Counted unit: the Lem whose source label is \texttt{key decomposition lemma} in the article (source0.tex lines 2078--2090, source label \texttt{key decomposition lemma}), delivered together with the article's own complete original proof of it (source0.tex lines 2091--2410, one \texttt{proof} environment of 320 lines). No mathematics is written, repaired or re-derived: statement and proof are the article's own text line for line. Required context retained uncounted: non-resonant Green's function (lines 932--941); Neumann series expansion (lines 945--948); difference in scales (lines 1601--1603); estimate on difference in scales (lines 1605--1608); random walk expansion (lines 1632--1635); length of in and out path (lines 1639--1641); estimate on difference in energy (lines 1829--1832); distance of lambda and E (lines 1878--1880); elements under eigenbasis (lines 1935--1937); define a(y) (lines 1967--1969); define vector bfa(y) (lines 1971--1973); transversality of norm of bfa (lines 1987--1989); eigenbasis block form (lines 2003--2012); eigenbasis schur complement (lines 2014--2016); Spectral stability lemma (lines 4489--4498). Every definition, display and label of those lines is retained unchanged. Omitted from this package and not redistributed: 1--931, 942--944, 949--1600, 1604--1604, 1609--1631, 1636--1638, 1642--1828, 1833--1877, 1881--1934, 1938--1966, 1970--1970, 1974--1986, 1990--2002, 2013--2013, 2017--2077, 2411--4488, 4499--4826. Nothing inside a retained excerpt is omitted or altered: every retained body is delivered exactly as the article's lines are written, so no omission receipt of any kind accompanies this package. Citations: the retained excerpts carry no citation command at all, so no bibliography entry of the article is transcribed and no reference is redistributed. Reference adaptation: none was needed and none was made; every reference inside the delivered unit resolves inside it, so no unresolved reference or citation is produced. Printed slips kept literal: no printed word, symbol, hypothesis or number of the counted statement or of its proof was altered. Layout and class: the article's own class is \texttt{amsart} with packages including amssymb, epsfig, booktabs, multirow, bbm, xltabular, tabularx, makecell, ragged2e, graphicx, color, amsmath, amsfonts, geometry; the wrapper is the lane's pinned \texttt{amsart} editorial wrapper at 10pt on A4, which restates the theorem-like environments the retained text needs and adds the article's own macro definitions that the retained text needs (\texttt{\textbackslash Z}, \texttt{\textbackslash barB}, \texttt{\textbackslash bfa}, \texttt{\textbackslash diffSj}, \texttt{\textbackslash dist}, \texttt{\textbackslash length}, \texttt{\textbackslash mcD}, \texttt{\textbackslash mcE}, \texttt{\textbackslash mcG}, \texttt{\textbackslash mcI}, \texttt{\textbackslash mcR}, \texttt{\textbackslash partialSj}, \texttt{\textbackslash partialSjminusone}, \texttt{\textbackslash smallSj}, \texttt{\textbackslash smallSjminusone}, \texttt{\textbackslash spec}, \texttt{\textbackslash tm}, \texttt{\textbackslash tq}, \texttt{\textbackslash tr}, \texttt{\textbackslash ts}, \texttt{\textbackslash tt}, \texttt{\textbackslash whn}, \texttt{\textbackslash wtGjminusone}, \texttt{\textbackslash wtSj}, \texttt{\textbackslash wtSjminusone}), with extra wrapper packages (bm, bbm, dsfont, xspace, cancel, mathtools). The delivered numbering is the unit's own. Assets: no retained span contains a figure environment, a graphics-inclusion command, a PDF-page inclusion, a table or a TikZ picture; no image file is copied into this package, the package declares no asset dependency and no image file is redistributed with it. Licence: version arXiv:2604.18989v1 of this article is distributed under the Creative Commons Attribution 4.0 International licence, as recorded in the arXiv licence metadata for this exact version and shown on the abstract page https://arxiv.org/abs/2604.18989v1. A full rights notice is delivered at licenses/arxiv-2604.18989-CC-BY-4.0.txt. Characters: every retained excerpt is pure ASCII, so the delivered engine, XeLaTeX with its default Latin Modern text font, reports no missing character.
\begin{lem}\label{non-resonant Green's function}
  If $\Lambda$ is non-resonant, then for any  $E\in [-\iota,4d+\beta+\iota]$, we have
  \begin{align}
   \label{non-resonant L2 norm} \|G_{\Lambda}(E)\|&\leq  \frac{1}{h-4d-\beta-\iota},\\
 % the off-diagonal decay
 % \begin{equation}
    \label{no-resonant off-diagonal decay} |G_{\Lambda}(x,y;E)| &\leq  \frac{1}{h-4d-\beta-\iota} \cdot \exp\{-\gamma_0 |x-y|_1\} \ {\rm for}\  \forall x,y \in \Lambda, 
  \end{align}
  where  $\gamma_0=\log(1+\frac{h-4d-\beta-\iota}{2d})>0$.
\end{lem}

  \begin{align}\label{Neumann series expansion}
  \notag  G_{\Lambda}(E) & = (2d-E+V_{\Lambda}-\Delta_{\Lambda})^{-1} \\
           &=(V_{\Lambda}+2d-E)^{-1} \sum_{s\geq 0} \left(  \Delta_{\Lambda} (V_{\Lambda}+2d-E)^{-1}\right). 
  \end{align}

\begin{equation}\label{difference in scales}
    \diffSj_\mcE := \wtSj_\mcE(B_j)-\wtSjminusone_\mcE(B_{j-1}).
\end{equation}

\begin{claim}\label{estimate on difference in scales}
    Assume $\mcE\in [-\iota,4d+\beta+\iota]$. Then 
  \[ \| \diffSj_\mcE  \| \leq g_{j-1}=\gamma^{(2a-0.3)L_{j-1}}.\]
\end{claim}

  \begin{align}\label{random walk expansion}
      % \diffSj_E (x,y) &= \sum_{w\in \partial^+ B_j \atop w\sim x} \sum_{x_1\in \partial^-\barB_{j-1}} \sum_{x'_1 \in \partial^+\barB_{j-1} \atop x'_1\sim x_1} \sum_{x'_2\in \partial^+\barB_{j-1}} \sum_{x_2\in \partial^- \barB_{j-1} \atop x_2\sim x'_2} \sum_{w'\in \partial^+ B_j \atop w'\sim y}
      \diffSj_\mcE (x,y) &= - \sum_{g\in \mcG_{\text{out,in}}(x,y) } G_{\barB_{j-1}\setminus B_{j}}(w,x_1;\mcE)   G_{\barB_j\setminus B_j}(x_1',x_2';\mcE)  G_{\barB_{j-1}\setminus B_{j}}(x_2,w';\mcE).  
  \end{align} 

  \begin{equation}\label{length of in and out path}
    \length(g)\geq 2\ \dist(\partial^+ B_j ,\partial^-\barB_{j-1})+4 \geq 2a L_{j-1}+2, 
  \end{equation}

\begin{claim}\label{estimate on difference in energy}
  Assume both $\mcE$ and $\mcE'$ are in $ [-\iota,4d+\beta+\iota]$. Then 
  \[ \| \partialSjminusone_{\mcE,\mcE'}  \| \leq \gamma^{1.9}|\mcE-\mcE'|.\]
\end{claim}

\begin{equation}\label{distance of lambda and E}
  \lambda \in \spec(\wtSjminusone_{\lambda}(B_{j-1})),\ |\lambda-E|<3g_{j-1}.
\end{equation}

\begin{equation}\label{elements under eigenbasis}
  \diffSj_{\lambda,m,m'}=-\left\langle \Gamma_{\partial \barB_{j-1}}  \wtGjminusone_{\lambda} \varphi_m,   G_{\barB_j\setminus B_j}(\lambda)  \Gamma_{\partial \barB_{j-1}} \wtGjminusone_{\lambda} \varphi_{m'} \right\rangle.
\end{equation}

\begin{equation}\label{define a(y)}
  a_m(y):= \Gamma_{\partial \barB_{j-1}}  \wtGjminusone_{\lambda} \varphi_m (y)
\end{equation}

\begin{equation}\label{define vector bfa(y)}
  \mathbf{a}(y)=(a_1(y),a_2(y),\cdots,a_{\whn}(y)).
\end{equation}

\begin{equation}\label{transversality of norm of bfa}
  |\bfa(\bar{y})| \geq \max_{y\in \partial ^+\barB_{j-1}} \mcI_{\psi_1} (y)  \geq \gamma^{(a+1+0.1)L_{j-1}}=t_{j-1}.
\end{equation}

\begin{equation}\label{eigenbasis block form}
\wtSjminusone_{\mcE}(B_{j-1}) = \begin{pmatrix}
q(\mcE) & r(\mcE) \\
s(\mcE) & t(\mcE)
\end{pmatrix}, \ 
\wtSj_{\mcE} (B_{j}) = \begin{pmatrix}
\tilde{q}(\mcE) & \tilde{r}(\mcE) \\
\tilde{s}(\mcE) & \tilde{t}(\mcE)
\end{pmatrix}.
\end{equation}

\begin{equation}\label{eigenbasis schur complement}
   \smallSjminusone_{\mcE}=q(\mcE)-r(\mcE)(t(\mcE)-\mcE)^{-1} s(\mcE) ,\ \smallSj_{\mcE}=\tilde{q}(\mcE)-\tilde{r}(\mcE)(\tilde{t}(\mcE)-\lambda)^{-1}\tilde{s}(\mcE).
\end{equation}

\begin{cor}[{\bf Spectral stability}]\label{Spectral stability lemma}
For each \(k = 1, \cdots, n\), we have
\[
\lambda_k(A + B) - \lambda_k(A) \in [\lambda_n(B), \lambda_1(B)]
\]
and thus 
\[
|\lambda_k(A + B) - \lambda_k(A)| \leq \|B\|.
\]
\end{cor}

\begin{lem}\label{key decomposition lemma}
   Assume that we choose the $a$-adic scale with $a\geq 5$. After conditioning on the randomness in $\barB_{j}\setminus \{\bar{y}\}$, we can decompose 
  {\small
  \begin{align*}
    \smallSj_E & = \frac{-1}{h+2d+\beta \omega(\bar{y})-\lambda}\left(\bfa(\bar{y})\otimes \bfa(\bar{y})+ \sum_{z\sim \bar{y} \atop z\in \partial^+\barB_{j-1}} \frac{1}{h+2d+\beta\omega(z)-\lambda } \cdot(\bfa(\bar{y}) \otimes  \bfa(z)+\bfa(z)\otimes\bfa(\bar{y})  )\right) \\
    \notag & \ \ \ + \mcD + \mcR(\omega(\bar{y})) .
  \end{align*}
  }
  Here $\mcD$ is a deterministic term, and the random term $\mcR(\omega(\bar{y}))$ satisfies
  \begin{equation}\label{estimate on mcR}
    \|\mcR(\omega(\bar{y}))\| \leq \frac{1}{2}\gamma^{2.5} |\bfa(\bar{y})|^2 + \frac{1}{2}\gamma^{0.5L_{j-1}} |\bfa(\bar{y})|^2.
  \end{equation}
\end{lem}

\begin{proof}[Proof of Lemma \ref{key decomposition lemma}]
 Consider 
 \begin{equation}\label{initial decomposition}
   \smallSj_E=\smallSjminusone_{\lambda}+(\smallSj_{\lambda}-\smallSjminusone_\lambda)+(\smallSj_E-\smallSj_{\lambda}).
 \end{equation}
 We analyze the above terms separately.\\

 \noindent {\Large   {\textbf{(Term $\smallSjminusone_{\lambda}$)}}  }\\
 
 \noindent We set 
 \begin{equation}\label{mcD1}
   \mcD_1=\smallSjminusone_{\lambda}.
 \end{equation}
 Obviously, $\mcD_1$ is deterministic, as it  only depends  on the randomness in $\barB_{j-1}$.\\




 \noindent {\Large {\textbf{(Term $(\smallSj_{\lambda}-\smallSjminusone_\lambda)$)}}     }\\

 \noindent By \eqref{eigenbasis block form} and \eqref{eigenbasis schur complement}, direct computation leads to 
\begin{equation}\label{baby decomposition}
  \smallSj_{\lambda}-\smallSjminusone_\lambda =(\tilde{q}(\lambda)-q(\lambda))- \tilde{r}(\lambda)(\tilde{t}(\lambda)-\lambda)^{-1}\tilde{s}(\lambda).
\end{equation}
Since $r(\lambda)=0$ and $s(\lambda)=0$, it is easy to see that  $\tq(\lambda)-q(\lambda),\tr(\lambda),\ts(\lambda)$ and $\tt(\lambda)-t(\lambda)$ contain elements of $\diffSj_{\lambda}$ under the eigenbasis, i.e.,  \eqref{elements under eigenbasis}.
Now under the notation \eqref{define a(y)} and \eqref{define vector bfa(y)}, \eqref{elements under eigenbasis} becomes
\begin{equation}\label{elements in bfa}
    \diffSj_{\lambda,m,m'}=-\sum_{x,y\in \partial^+ \barB_{j-1}}a_{m}(x)G_{\barB_j\setminus B_j }(x,y;\lambda) a_{m'}(y).
\end{equation}
We expand $G_{\barB_{j}\setminus B_j}(\lambda)$ via the Neumann series similar to  \eqref{Neumann series expansion}, leading to the following random walk expansion:
\begin{align}\label{K random walk expansion}
  G_{\barB_{j}\setminus B_j}(x,y;\lambda) =\sum_{g\in \mcG_{\partial^+ \barB_{j-1}}(x,y)}  \frac{1}{V_{x}+2d-\lambda}\cdot \frac{1}{V_{x_1}+2d-\lambda}\cdots \frac{1}{V_{x_{s-1}}+2d-\lambda}\cdot \frac{1}{V_{y}+2d-\lambda}.
\end{align}
Here for $x,y\in \partial^+ \barB_{j-1}$, the set $\mcG_{\partial^+ \barB_{j-1}}(x,y)$ contains all random walks  $g$ contained in $\barB_{j}\setminus B_j$ with 
\[g:\ x\sim x_1\sim x_2\cdots \sim x_{s-1}\sim y,\ s=\length(g) .\]
We decompose $\mcG_{\partial^+ \barB_{j-1}}(x,y)$ into disjoint unions 
\[\mcG_{\partial^+ \barB_{j-1}}(x,y) =\mcG_{\partial^+ \barB_{j-1}}^{(0)} (x,y)\cup \mcG_{\partial^+ \barB_{j-1}}^{(1)} (x,y)\cup \mcG_{\partial^+ \barB_{j-1}}^{(2)} (x,y) ,\]
with
\[\mcG_{\partial^+ \barB_{j-1}}^{(0)} (x,y)= \left\{ g\in \mcG_{\partial^+ \barB_{j-1}} (x,y) :\  \bar{y}\notin g \right\},\]
\[\mcG_{\partial^+ \barB_{j-1}}^{(1)} (x,y)= \left\{ g\in \mcG_{\partial^+ \barB_{j-1}} (x,y) : \ \bar{y}\in g ,\ \length(g)=0 \ {\rm or} \ 1 \right\},\]
\[\mcG_{\partial^+ \barB_{j-1}}^{(2)} (x,y)= \left\{ g\in \mcG_{\partial^+ \barB_{j-1}} (x,y) : \ \bar{y}\in g ,\ \length(g)\geq 2 \right\}.\]
Clearly, $\mcG_{\partial^+ \barB_{j-1}}^{(1)}$ is not empty if and only if $x=\bar{y}$ or $y=\bar{y}$. Therefore, \eqref{K random walk expansion} can be decomposed into 
\begin{equation}\label{decomposition of K}
  K(x,y):= G_{\barB_{j}\setminus B_j}(x,y;\lambda) =K_0(x,y)+K_1(x,y)+K_2(x,y),
\end{equation}
where 
\begin{equation}\label{Ki comes from graph decomposition}
  K_i(x,y)= \sum_{g\in \mcG_{\partial^+ \barB_{j-1}}(x,y)}  \frac{1}{V_{x}+2d-\lambda}\cdot \frac{1}{V_{x_1}+2d-\lambda}\cdots \frac{1}{V_{y}+2d-\lambda}, \ i=1,2,3 .
\end{equation}
Clearly, our decomposition ensures that $K_1(x,y),K_2(x,y)$ depend only  on $V_{\text{r}}(\bar{y})=\omega(\bar{y})$, whereas $K_0(x,y)$ does not. Since $\barB_{j}\setminus B_j$ is non-resonant, we have 
\begin{equation}\label{K1 formula 1}
  K_1(\bar{y},\bar{y})= \frac{1}{h+2d+\beta \omega(\bar{y})-\lambda} 
\end{equation}
and
\begin{equation}\label{K1 formula 2}
  K_1(\bar{y},z)=K_1(z,\bar{y})=\frac{1}{h+2d+\beta \omega(\bar{y})-\lambda} \cdot \frac{1}{h+2d+\beta \omega(z)-\lambda} \ {\rm for}\ \forall z\sim \bar{y},z\in \partial^+\barB_{j-1}.
\end{equation} 
Moreover, the random walks in $\mcG_{\partial^+ \barB_{j-1}}^{(2)}$ have length larger than $2$.
We use $\lambda\in [-\iota,4d+\beta+\iota]$ to estimate
\begin{align}\label{K2 estimate}
\notag  |K_2(x,y)| & \leq \sum_{t \geq (|x-\bar{y}|_1+|\bar{y}-y|_1)\vee 2 \atop r+s=t,r>|x-\bar{y}|_1,s>|\bar{y}-y|_1} \sum_{ g_1 =(x_0=x, x_1,\cdots,x_r=\bar{y}) \atop g_1\in \barB_j\setminus B_j} \sum_{ g_2 =(x_r=\bar{y}, x_{r+1},\cdots,x_r=\bar{y}) \atop g_2\in \barB_j\setminus B_j} \frac{1}{V_{x}+2d-\lambda}\cdots \frac{1}{V_{y}+2d-\lambda}  \\
 \notag  &\leq \sum_{t \geq (|x-\bar{y}|_1+|\bar{y}-y|_1)\vee 2 } (t+1) (2d)^t \left(\frac{1}{h-2d-\beta-\iota}\right)^{t+1} \\
\notag & \lesssim_d \left(\frac{2d}{h-2d-\beta-\iota}\right)^{(|x-\bar{y}|_1+|\bar{y}-y|_1)\vee 2 +1}\\
      &<\gamma^{0.9[(|x-\bar{y}|_1+|\bar{y}-y|_1)\vee 3]}.
\end{align}
Here the validity of \eqref{K2 estimate} also requires 
\begin{equation}\label{essential barrier high 3}
   \frac{2d}{h-2d-\beta-\iota}\leq \frac{2d}{h-2d-2}\lesssim_d   \left(\frac{1}{4d+h+1}\right) ^{1.9} \leq  \gamma^{0.9},
\end{equation}
which will hold when $h$ is sufficiently large (depending only on $d$). 
We first discuss the term $\tq(\lambda)-q(\lambda)$ in \eqref{baby decomposition}. For $m,m'\in \{1,\cdots,\whn \}$, using \eqref{elements in bfa} and the decomposition \eqref{decomposition of K} yields 
\begin{align*}
  (\tq(\lambda)-q(\lambda))_{m,m'} & = \diffSj_{\lambda,m,m'}= -\sum_{x,y\in \partial^+ \barB_{j-1}}a_{m}(x) K(x,y) a_{m'}(y)\\
    &=-\sum_{x,y\in \partial^+ \barB_{j-1}}a_{m}(x)\left(K_0(x,y) +K_1(x,y)+K_2(x,y) \right)a_{m'}(y).
\end{align*}
We set 
\begin{equation}\label{mcD2}
  \mcD_{2,m,m'}=-\sum_{x,y\in \partial^+ \barB_{j-1}}a_{m}(x)K_0(x,y)a_{m'}(y),
\end{equation}
which is deterministic since $\bfa(x),\bfa(y)$ and $K_0(x,y)$ are independent of $\omega(\bar{y})$.
Moreover, recalling t \eqref{K1 formula 1} and \eqref{K1 formula 2}, a direct computation yields 
{\small
\begin{align*}
  &\sum_{x,y\in \partial^+ \barB_{j-1}}a_{m}(x)K_1(x,y)a_{m'}(y) \\
  =& \frac{1}{h+2d+\beta\omega(\bar{y})-\lambda}a_{m}(\bar{y}) a_{m'}(\bar{y}) + \sum_{z\sim \bar{y} \atop z\in \partial^+\barB_{j-1}} \frac{a_m(\bar{y})a_{m'}(z)+a_m(z)a_{m'}(\bar{y})}{(h+2d+\beta\omega(\bar{y})-\lambda)\cdot (h+2d+\beta\omega(z)-\lambda)} \\
  =&\frac{1}{h+2d+\beta \omega(\bar{y})-\lambda}\left(\bfa(\bar{y})\otimes \bfa(\bar{y})+ \sum_{z\sim \bar{y} \atop z\in \partial^+\barB_{j-1}} \frac{1}{h+2d+\beta\omega(z)-\lambda } \cdot(\bfa(\bar{y}) \otimes  \bfa(z)+\bfa(z)\otimes\bfa(\bar{y})  )\right)_{m,m'}.
\end{align*}
}
Next, we set
\begin{equation}\label{mcR1}
  \mcR_{1,m,m'}(\omega(\bar{y}))=-\sum_{x,y\in \partial^+ \barB_{j-1}}a_{m}(x)K_2(x,y)a_{m'}(y).
\end{equation}
\begin{claim}\label{estimate on mcR1}
  We have 
  \[\|\mcR_{1}(\omega(\bar{y}))\| \leq \frac{1}{2}\gamma^{2.5} |\bfa(\bar{y})|^2. \]
\end{claim}
\begin{proof}[Proof of Claim \ref{estimate on mcR1}]
  By our definition,
  \[\mcR_1(\omega(\bar{y}))=-\sum_{x,y\in \partial^+ \barB_{j-1}}K_2(x,y)\cdot \bfa(x)\otimes \bfa(y).\]
  Hence, using the Minkowski inequality and \eqref{K2 estimate} yields
  \begin{align*}
    \| \mcR_{1}(\omega(\bar{y})) \| & \leq \sum_{x,y\in \partial^+ \barB_{j-1}} |K_2(x,y)|\cdot \| \bfa(x)\otimes \bfa(y) \| = \sum_{x,y\in \partial^+ \barB_{j-1}} |K_2(x,y)|\cdot | \bfa(x)|\cdot |\bfa(y) | \\
          &\leq |\bfa(\bar{y})|^2 \sum_{x,y\in \partial^+ \barB_{j-1}} \gamma^{0.9[(|x-\bar{y}|_1+|\bar{y}-y|_1)\vee 3]} \leq |\bfa(\bar{y})|^2 \sum_{x,y\in \Z^d } \gamma^{0.9[(|x-\bar{y}|_1+|\bar{y}-y|_1)\vee 3]} \\
          &\leq |\bfa(\bar{y})|^2 \sum_{r,s\geq 0} (2d)^{r+s}\gamma^{0.9[(r+s)\vee 3]} \lesssim_d |\bfa(\bar{y})|^2 \gamma^{0.9\times 3} \\
          &\leq \frac{1}{2}\gamma^{2.5} |\bfa(\bar{y})|^2 
  \end{align*}
  supposing $h$ is large enough (depending only on $d$) to make sure  $0<\gamma\ll1$.  In the above argument, we also use the fact that $\bar{y}$ is chosen to maximize $|\bfa(y)|$.

\end{proof}
Summarize the above discussions, we have decomposed
{\small
\begin{align}\label{further baby decomposition 1}
  \tq(\lambda)-q(\lambda)&= \mcD_2 +\mcR_1 (\omega(\bar{y}))\\
  \notag &\ \ \ -\frac{1}{h+2d+\beta \omega(\bar{y})-\lambda}\left(\bfa(\bar{y})\otimes \bfa(\bar{y})+ \sum_{z\sim \bar{y} \atop z\in \partial^+\barB_{j-1}} \frac{1}{h+2d+\beta\omega(z)-\lambda } \cdot(\bfa(\bar{y}) \otimes  \bfa(z)+\bfa(z)\otimes\bfa(\bar{y})  )\right).
\end{align}
}

For the remaining term $\tilde{r}(\lambda)(\tilde{t}(\lambda)-\lambda)^{-1}\tilde{s}(\lambda)$ in \eqref{baby decomposition}, we denote it by $\mcR_2(\omega(\bar{y}))$. It can be estimated  as follows:
\begin{claim}\label{estimate on mcR2}
  For 
  \begin{equation}\label{mcR2}
    \mcR_2(\omega(\bar{y}))=-\tilde{r}(\lambda)(\tilde{t}(\lambda)-\lambda)^{-1}\tilde{s}(\lambda),
  \end{equation}
  we have 
  \[ \|\mcR_2(\omega(\bar{y})) \|\leq \frac{1}{6}\gamma^{(a-2)L_{j-1}} |\bfa(\bar{y})|^2 .\] 
\end{claim}
\begin{proof}[Proof of Claim \ref{estimate on mcR2}]
  For $m,m'\in \{1,\cdots,\whn\}$, we expand elements in $\mcR_2(\omega(\bar{y}))$ into 
  \begin{align*}
    \mcR_{2,m,m'}(\omega(\bar{y}))=-\sum_{ \whn+1\leq \tm,\tm'\leq  \# B_j } \diffSj_{\lambda,m,\tm}\cdot [(\tt(\lambda)-\lambda)^{-1}]_{\tm,\tm'}\cdot \diffSj_{\lambda,\tm',m'}.
  \end{align*}
  By \eqref{elements in bfa}, this can be further expanded into 
  \begin{align}\label{mcR2 expansion}
        &\mcR_{2,m,m'}(\omega(\bar{y}))\\
     \notag   &=-\sum_{x,x',y,y'\in \partial^+\barB_{j-1}} \sum_{ \whn+1\leq \tm,\tm'\leq  \# B_j } a_{m}(x)K(x,x') a_{\tm}(x')\cdot [(\tt(\lambda)-\lambda)^{-1}]_{\tm,\tm'}\cdot a_{\tm'}(y')K(y',y) a_{m'}(y).
  \end{align}
  Now, recall that we choose the eigenbasis such that $\mcE_m\notin I_{\frac{\varepsilon_{j-1}}{2}}(E)$ for $\whn\leq m\leq \# B_j$, and we have \eqref{distance of lambda and E}. Thus 
  \[\dist(\spec(t(\lambda)),\lambda)> \frac{\varepsilon_{j-1}}{2}-3g_{j-1}.\]
  Moreover, since $\tt(\lambda)-t(\lambda)$ is the restriction of $\diffSj_{\lambda}$ on the subspace spanned by ${\varphi_{\whn+1},\cdots,\varphi_{\#B_j}}$, applying Claim \ref{estimate on difference in scales} yields 
  \[\| \tt(\lambda)-t(\lambda) \| \leq \|\diffSj_{\lambda}\|\leq g_{j-1}.\]
  Therefore, using  Corollary \ref{Spectral stability lemma},  we have 
  \begin{align*}
      \dist(\spec(\tt(\lambda)),\lambda) & >\dist(\spec(t(\lambda)),\lambda)-g_{j-1}>\frac{\varepsilon_{j-1}}{2}-4g_{j-1} \\
        &=\frac{1}{2}\gamma^{(a+1+0.2)L_{j-1}}-4\gamma^{(2a-0.3)L_{j-1}} >\frac{\varepsilon_{j-1}}{3},
  \end{align*}
  which implies 
  \begin{equation}\label{resolvent of tilde t}
    \|(\tt(\lambda)-\lambda)^{-1}\|< \frac{3}{\varepsilon_{j-1}}.
  \end{equation}

  On the other hand, the definition \eqref{define a(y)} ensures that 
  \begin{align*}
    |a_m(y)| & = |\Gamma_{\partial \barB_{j-1}}  \wtGjminusone_{\lambda} \varphi_m (y)| =| \Gamma_{\partial \barB_{j-1}} G_{\barB_{j-1}\setminus B_{j-1}}(\lambda)  \Gamma_{j-1} \varphi_m(y)| \\
       & \leq \sum_{w\in \partial^{-}\barB_{j-1} \atop  w\sim y} \sum_{w'\in \partial^+ B_{j-1} \atop y'\in \partial^-B_{j-1} , w' \sim y'}  |G_{\barB_{j-1}\setminus B_{j-1}}(w,w';\lambda)|\cdot |\varphi_m(y')| \\ 
       &\leq \#\partial^+ B_{j-1} \cdot \frac{(2d)^2}{h-4d-\beta-\iota} \left(\frac{2d}{h-2d-\beta-\iota}\right)^{\dist(\partial^+B_{j-1},\partial^-\barB_{j-1})}.
  \end{align*}
  Here we used Lemma \ref{non-resonant Green's function} to estimate $G_{\barB_{j-1}\setminus B_{j-1}}(\lambda)$. Now by the fact that  
  \[\dist(\partial^+B_{j-1},\partial^-\barB_{j-1}) =aL_{j-1}-1\]
  and taking $h$ large enough (depending on $d,\beta$), we obtain 
  \begin{equation}\label{decay of a(y)}
    |a_m(y)|<\gamma^{(a-0.1)L_{j-1}}\  {\rm for}\ \forall 1\leq m\leq \#B_j,\  \forall y\in \partial^+\barB_{j-1}.
  \end{equation}

  We can use \eqref{resolvent of tilde t} and \eqref{decay of a(y)} to bound 
  \begin{align*}
    \sum_{ \whn+1\leq \tm,\tm'\leq  \# B_j }  \bigg|a_{\tm}(x')\cdot [(\tt(\lambda)-\lambda)^{-1}]_{\tm,\tm'}\cdot a_{\tm'}(y')\bigg|& \leq  (\# B_j)^2 \gamma^{2(a-0.1)L_{j-1}} \cdot \frac{3}{\gamma^{(a+1+0.2)L_{j-1}}} \\
      &< \gamma^{(a-1-0.5)L_{j-1}}.
  \end{align*}
  Plugging this estimate into \eqref{mcR2 expansion} gives  
  \begin{align*}
    |\mcR_{2,m,m'}(\omega(\bar{y}))| &\leq\gamma^{(a-1-0.5)L_{j-1}} \bigg| \sum_{x,x',y,y'\in \partial^+\barB_{j-1}} a_{m}(x)K(x,x')  K(y',y) a_{m'}(y) \bigg| \\
    &\leq \gamma^{(a-1-0.5)L_{j-1}} |\bfa(\bar{y})|^2 \cdot (\# \partial^+\barB_{j-1})^4 \|K\|^2.
  \end{align*}
  In the above argument,  we again used the fact that $\bar{y}$ maximizes $|\bfa(y)|$. Finally, recall that $K$ is the Green's function $G_{\barB_{j}\setminus B_j}(x,y;\lambda)$. Applying Claim \ref{non-resonant Green's function} and the Schur's test yields
  \begin{align*}
    \|\mcR_{2}(\omega(\bar{y}))\| & \leq \frac{\gamma^{(a-1-0.5)L_{j-1}} }{(h-4d-\beta-\iota)^2}|\bfa(\bar{y})|^2 \cdot (\# \partial^+\barB_{j-1})^4 \cdot (\# B_{j})^2 \\
       &\ll \frac{1}{6}\gamma^{(a-2)L_{j-1}} |\bfa(\bar{y})|^2.
  \end{align*}
The absorption of terms such as $\# B_{j}$ is possible because we can take $k$ sufficiently large, thereby making $L_{j-1}$ large enough.
\end{proof}


It follows from \eqref{baby decomposition}, \eqref{further baby decomposition 1} and \eqref{mcR2} that
{\small
\begin{align}\label{second term decomposition}
  \smallSj_{\lambda}-\smallSjminusone_\lambda &= \mcD_2 +\mcR_1 (\omega(\bar{y})) +\mcR_2(\omega(\bar{y})) \\
  \notag &\ \ \ -\frac{1}{h+2d+\beta \omega(\bar{y})-\lambda}\left(\bfa(\bar{y})\otimes \bfa(\bar{y})+ \sum_{z\sim \bar{y} \atop z\in \partial^+\barB_{j-1}} \frac{1}{h+2d+\beta\omega(z)-\lambda } \cdot(\bfa(\bar{y}) \otimes  \bfa(z)+\bfa(z)\otimes\bfa(\bar{y})  )\right).
\end{align}
}


\noindent {\Large  {\textbf{(Term $(\smallSj_E-\smallSj_{\lambda})$)}}  }\\

\noindent For this term, we compute it as 
\begin{align}\label{baby decomposition 2}
 \notag \smallSj_E-\smallSj_{\lambda} &=(\tq(E)-\tq(\lambda)) +\tr(\lambda) (\tt(\lambda)-\lambda)^{-1} \ts(\lambda)  -\tr(E) (\tt(E)-E)^{-1} \ts(E) \\
                     & =(q(E)-q(\lambda)) \\
         \notag                      & \ \ \ +[(\tq(E)-q(E))-(\tq(\lambda) -q(\lambda)) ]\\
         \notag                     &\ \ \ +\tr(\lambda) (\tt(\lambda)-\lambda)^{-1} \ts(\lambda)  -\tr(E) (\tt(E)-E)^{-1} \ts(E).  
\end{align}
We set 
\begin{equation}\label{mcD3}
  \mcD_3= (q(E)-q(\lambda)),
\end{equation}
which is deterministic, since by our construction,  $q(\mcE)$ and $\lambda$ are independent of $\omega(\bar{y})$. 

Next, for the second line in \eqref{baby decomposition 2}, we set 
\begin{equation}\label{mcR3}
  \mcR_3(\omega(\bar{y}))= (\tq(E)-q(E))-(\tq(\lambda) -q(\lambda)).
\end{equation}
\begin{claim}\label{estimate on mcR3}
  We have 
  \[\|\mcR_3(\omega(\bar{y})) \| \leq \frac{1}{6}\gamma^{(a-2)L_{j-1}} |\bfa(\bar{y})|^2 .\]
\end{claim}
\begin{proof}[Proof of Claim \ref{estimate on mcR3}]
Indeed, $\mcR_3(\omega(\bar{y}))$ can be viewed  as the matrix 
\[\diffSj_{E}-\diffSj_{\lambda} \]
restricted to the subspace spanned by $\{\varphi_1,\cdots,\varphi_{\whn}\}$. Therefore,
\[\|\mcR_3(\omega(\bar{y})) \|\leq \| \diffSj_{E}-\diffSj_{\lambda}  \| .\]
Expanding the Green's function in \eqref{random walk expansion} further by a Neumann series yields 
\begin{equation}\label{scales not derivation}
  \diffSj_{\mcE}(x,y) =-\sum_{g\in \mcG_{\text{out,in}}(x,y) } \frac{1}{V_{x_1}+2d-\mcE} \cdot \frac{1}{V_{x_2}+2d-\mcE} \cdots \frac{1}{V_{x_{s-1}}+2d-\mcE} 
\end{equation}
with the random walk 
\[g=(x_0=x,x_1,x_2,\cdots,x_{s-1},x_s=y),x_{i-1}\sim x_i\]
in $\mcG_{\text{out,in}}(x,y)$. (For the definition of $\mcG_{\text{out,in}}(x,y)$, see the proof of Claim \ref{difference in scales}.) 

Taking the  derivative  about $\mcE$ on \eqref{scales not derivation} and applying \eqref{length of in and out path} yield  that 
\begin{align*}
  |\frac{d}{d \mcE} (\diffSj_{\mcE}(x,y))| & \leq \sum_{g\in \mcG_{\text{out,in}}(x,y) } \left(\frac{1}{h-2d-\beta-\iota}\right) ^{\length(g)} \\
     &\leq \sum_{s\geq 2aL_{j-1}+2}  \left(\frac{2d}{h-2d-\beta-\iota}\right) ^s <\gamma^{(2a-0.1)L_{j-1}}
\end{align*}
holds for all $\mcE\in [-\iota,4d+\beta+\iota]$ and $x,y\in \partial^-B_j$,  provided 
\[ \left(\frac{2d}{h-2d-\beta-\iota}\right)^{2a}< \gamma^{(2a-0.1)}. \]
This can be ensured by 
\begin{equation}\label{essential barrier high 4}
  \left(\frac{2d}{h-2d-2}\right)^{2a}< \left(\frac{1}{4d+h+1}\right)^{(2a-0.1)},
\end{equation}
which requires $h$ to be sufficiently large (depending only on $d$). Thus by the mean value theorem and \eqref{distance of lambda and E},  
\[ |\diffSj_{E}(x,y)-\diffSj_{\lambda}(x,y)| < \gamma^{(2a-0.1)L_{j-1}} |\lambda-E| < 3 g_{j-1} \gamma^{(2a-0.1)L_{j-1}} .\]
Apply the Schur's  test, we get  
\begin{align*}
  \| \diffSj_E-\diffSj_{\lambda}  \|& \leq (\# \partial^-B_j)^2 \cdot 3 g_{j-1} \gamma^{(2a-0.1)L_{j-1}} \leq   \gamma^{(4a-0.5)L_{j-1}} \\
     & = \gamma^{(2a-2-0.7)L_{j-1}} t^2_{j-1} \leq \frac{1}{6} \gamma^{(a-2)L_{j-1}} |\bfa(\bar{y})|^2.
\end{align*}
In the above estimate, we used \eqref{transversality of norm of bfa}. Therefore, we  have proven   
\begin{equation}\label{bound on first part of mcR3}
  \|\mcR_3(\omega(\bar{y})) \| \leq \frac{1}{6}\gamma^{(a-2)L_{j-1}} |\bfa(\bar{y})|^2.
\end{equation}
\end{proof}


Finally, for the third line in \eqref{baby decomposition 2}, we set 
\begin{equation}\label{mcR4}
  \mcR_4(\omega(\bar{y})) =\tr(\lambda) (\tt(\lambda)-\lambda)^{-1} \ts(\lambda)  -\tr(E) (\tt(E)-E)^{-1} \ts(E).
\end{equation}
\begin{claim}\label{estimate on mcR4}
  Assume that we choose the $a$-adic scale with $a\geq 5$. Then we have
  \[\|  \mcR_4(\omega(\bar{y}))  \| \leq \frac{1}{6}\gamma^{0.5L_{j-1}} |\bfa(\bar{y})|^2 .\]
\end{claim}
\begin{proof}[Proof of Claim \ref{estimate on mcR4}]
  Compute $\mcR_4(\omega(\bar{y}))$ as follows:
  \begin{align*}
    \mcR_4(\omega(\bar{y})) &=        (\tr(\lambda) -\tr(E)) (\tt(\lambda)-\lambda)^{-1} \ts(\lambda) \\
             \notag               &\qquad + \tr(E)(\tt(\lambda)-\lambda)^{-1} (\ts(\lambda)-\ts(E)) \\
                   \notag         &\qquad + \tr(E) (\tt(\lambda)-\lambda)^{-1} \cdot [(\tt(E)-\tt(\lambda))-(E-\lambda) ]\cdot (\tt(E)-E)^{-1} \ts(E).
  \end{align*}

We have the following  available estimates:
\begin{itemize} 
 \item By $r(\lambda)=0,s(\lambda)=0$, $|\lambda-E|<3g_{j-1}$ and Claim \ref{estimate on difference in energy}, we have 
             \[\| s(E) \|,\ \| r(E)\| \leq \| \partialSjminusone_{\lambda,E}\|\leq \gamma^{1.9} \cdot 3g_{j-1} \]
             and 
           \[\| \ts(\lambda)-\ts(E) \|,\ \| \tr(\lambda)-\tr(E)\| \leq \| \partialSj_{\lambda,E}\|\leq \gamma^{1.9} \cdot 3g_{j-1} .\]
  \item By $r(\lambda)=0,s(\lambda)=0$ and Claim \ref{difference in scales}, we have 
               \[\|\tr(\lambda)\|, \|\ts(\lambda)\|,\   \| \tr(E)-r(E)\|, \ \| \ts(E)-s(E) \|,\  \|\tt(E)-t(E)\| \leq g_{j-1}.\] 
       Therefore, we also have 
       \[  \| \ts(E) \|,\ \| \tr(E)\| \leq g_{j-1}+ \gamma^{1.9} \cdot 3g_{j-1}\leq 2 g_{j-1}.   \]
  \item By similar argument leading  to \eqref{resolvent of tilde t}, and the fact that 
  %  \[|\lambda-\wtlambda|<9g_{j-1},\]
      \[ \| \tt(\lambda)-\tt(E)\| \leq \| \partialSj_{\lambda,E}\|\leq \gamma^{1.9} g_{j-1},\]
        we have  
        \[\|(\tt(\lambda)-\lambda)^{-1} \|, \ \|(\tt(E)-E)^{-1} \| <\frac{4}{\varepsilon_{j-1}}. \]
\end{itemize}
Use the estimates listed above, we can control  the norm of $\mcR_4(\omega(\bar{y}))$ as 
\begin{align}\label{bound on mcA}
  \notag \|\mcR_4(\omega(\bar{y}))\| &\lesssim \frac{g_{j-1}^2 }{\varepsilon_{j-1}} + \frac{g_{j-1}^3 }{\varepsilon_{j-1}^2} = \gamma^{(3a-1-0.8)L_{j-1}} +\gamma^{(4a-2-1.4)L_{j-1}}\\
    &\lesssim \gamma^{[3a-1-0.8-2(a+1+0.1)]L_{j-1}}t^2_{j-1} \leq \frac{1}{6}\gamma^{0.5L_{j-1}} t^2_{j-1}.
\end{align}
The validity of \eqref{bound on mcA} requires 
\begin{equation}\label{range of a}
  3a-1-0.8-2(a+1+0.1)>0.5,
\end{equation}
which means we need to take $a\geq 5$. Thus, using \eqref{transversality of norm of bfa} and \eqref{bound on mcA} yields 
\[\| \mcR_4(\omega(\bar{y})) \|\leq \frac{1}{6}\gamma^{0.5L_{j-1}} |\bfa(\bar{y}) |^2 .\]

\end{proof}
%\begin{rmk}
%From \eqref{range of a}, treating $0.1$ as a small quantity, the parameter $a$ must satisfy $a>3+$. Consequently, the dyadic scale ($a=2$) in \cite{Imb21} is inadequate here.
%\end{rmk}



Combining \eqref{baby decomposition 2}, \eqref{mcD3}, \eqref{mcR3} and \eqref{mcR4} together  yields the decomposition
\begin{equation}\label{third term decomposition}
  \smallSj_E-\smallSj_{\lambda} = \mcD_3+\mcR_3(\omega(\bar{y}))+\mcR_4(\omega(\bar{y})).
\end{equation}





 \noindent {\Large  {\textbf{(Summarize all terms)}}     }\\

\noindent Finally, by \eqref{initial decomposition}, \eqref{mcD1}, \eqref{second term decomposition} and \eqref{third term decomposition}, we take 
\[\mcD=\mcD_1+\mcD_2+\mcD_3\]
and 
\[\mcR(\omega(\bar{y}))= \mcR_1(\omega(\bar{y}))+\mcR_2(\omega(\bar{y}))+\mcR_3(\omega(\bar{y}))+\mcR_4(\omega(\bar{y})).\]
Then the decomposition in Lemma \ref{key decomposition lemma} is obtained. Using Claim \ref{estimate on mcR1}, Claim \ref{estimate on mcR2}, Claim \ref{estimate on mcR3} and Claim \ref{estimate on mcR4} on the random parts  yield \eqref{estimate on mcR}. We have finished  the proof of Lemma \ref{key decomposition lemma}.
\end{proof}

\end{document}
